\documentclass[10pt,a4paper]{amsart}
\usepackage[margin=19mm]{geometry}
\usepackage{amsmath,amssymb,mathtools}
\usepackage[scr=rsfs,cal=euler]{mathalfa}
\usepackage{xfrac}
\usepackage[unicode,hidelinks,bookmarks=false]{hyperref}
\newcommand{\ie}{i.e.\ }
\newcommand{\cf}{cf.\ }
\newcommand{\resp}{respectively\ }
\newcommand{\Subsection}{\subsection}
\providecommand{\nobreakdash}{\nobreak\hbox{-}}
\setlength{\emergencystretch}{2em}
\multlinegap0pt
\usepackage{amssymb}
\newtheorem{lemma}{Lemma}
\newtheorem{theorem}{Theorem}
\title{On binary codes that are maximal\\ totally isotropic subspaces\\ with respect to an alternating form}
\author{Patrick King, Mikhail Kochetov}
\date{Source: arXiv:2510.05464v1 (CC BY 4.0)}
\begin{document}
\maketitle

\noindent\emph{Source and licence.} On binary codes that are maximal\\ totally isotropic subspaces\\ with respect to an alternating form. Version arXiv:2510.05464v1 (CC BY 4.0), distributed by arXiv under the Creative Commons Attribution 4.0 International licence, http://creativecommons.org/licenses/by/4.0/, as recorded in the arXiv licence metadata for this exact version and shown on the abstract page https://arxiv.org/abs/2510.05464v1. Full licence notice: \texttt{licenses/arxiv-2510.05464-CC-BY-4.0.txt}.\par\smallskip

\noindent\textit{Editorial scope.} Counted unit: the article's theorem W\_L:n=0(8) (source0.tex lines 1170--1215). The article's complete original proof of it is delivered with it (source0.tex lines 1217--1305, one \texttt{proof} environment of 89 lines). No mathematics is written, repaired or re-derived: the statement and the proof are the article's own lines, in the article's own notation, and no retained line is substituted or re-flowed.  Retained uncounted as the source-local context the proof needs: lines 871--888; lines 1160--1162.  Nothing was omitted from any retained span, so the package carries no omission record.  No retained line contains a citation, so the wrapper carries no bibliography and no bibliography entry.  No retained line contains a figure environment, an image-inclusion command or drawing code, so no image is delivered and the package declares no asset dependency.  Editorial formatting only, in addition: an AMS \texttt{amsart} preamble that restates the article's own declarations and macros used by the retained lines, namely the environments \texttt{lemma}, \texttt{theorem} on separate counters, together with the standard packages those lines need.  The retained environments print in this document as Lemma 1, Theorem 1 in order of appearance, whereas the article prints the same items with the numbering of its own full document.  The complete native source of the article as distributed in the arXiv e-print for this exact version is bundled unchanged as \texttt{sources/186531/source0.tex}.
\begin{lemma}\label{semi-invGI}
    The non-trivial \(1\)-dimensional characters of \(G\) have associated semi-invariants described by the following table:
    \begin{center}
        \begin{tabular}{|c|c|c|c|}
            \hline
             Character & Value on \(A\) & Value on \(B\) & Semi-invariant \\ \hline
             \(\chi_1\) & \(-1\) & \(-i\) & \(p = 2xy(y^2-x^2)(y^2+x^2)\) \\
             \(\chi_2\) & \(1\) & \(-1\) & \(w = x^2y^2(x^4-y^4)^2\) \\ 
             \(\chi_3\) & \(-1\) & \(1\) & \(u = -\frac{1}{2}(x^{12} - 33x^8y^4 - 33x^4y^8 + y^{12})\) \\ 
             \(\chi_4\) & \(1\) & \(-i\) & \(D_1 = pu\) \\ 
             \(\chi_5\) & \(1\) & \(i\) & \(D_2 = puw\) \\ 
             \(\chi_6\) & \(-1\) & \(i\) & \(D_3 = pw\) \\ 
             \(\chi_7\) & \(-1\) & \(-1\) & \(D_4 = uw\) \\ 
             \hline
        \end{tabular}
    \end{center}
    Moreover, each of the listed semi-invariants generates its corresponding \(\mathbb{C}[s,t]\)-module of semi-invariants.
\end{lemma}

\begin{lemma}\label{repG_n=0(4)}
    Set \(e_1 = W_L^+(x,y),\, e_2 = W_L^-(x,y),\, e_3 = W_L^+(y,x),\, e_4 = W_L^-(y,x)\) and suppose that \(n \equiv 0 \ (\textrm{mod}\ 4)\). Then, \(V\) is \(G\)-invariant and its decompositions into irreducible \(G\)-submodules is given by \(V = \langle v_0 \rangle \oplus \langle v_1, v_2, v_3 \rangle\) where \(v_0 = e_1 + e_3,\, v_1 = e_2 + e_4,\, v_2 = e_1 - e_3\), and \(v_3 = e_2 - e_4\).
\end{lemma}

\begin{theorem}\label{W_L:n=0(8)}
    Suppose that \(n \equiv 0 \ (\mathrm{mod}\ 4)\) and that \(L \subset \mathbb{F}_2^n\) is a Lagrangian subspace of Type II. Set \(v_0 = W_L^+(x,y) + W_L^+(y,x),\, v_1 = W_L^-(x,y) + W_L^-(y,x),\, v_2 = W_L^+(x,y) - W_L^+(y,x)\), and \(v_3 = W_L^-(x,y) - W_L^-(y,x)\). Then, \(n \equiv 0 \ (\mathrm{mod}\ 8)\) and
    \begin{align*}
        &v_0 \in \mathbb{C}[s,t] \\
        &v_1 \in p(q_1r_3+q_3r_1)\mathbb{C}[s,t] \oplus (p_1r_3 - p_3r_1)\mathbb{C}[s,t] \oplus -pw(p_1q_3+p_3q_1)\mathbb{C}[s,t] \\
        &v_2 \in p(q_3r_2-q_2r_3)\mathbb{C}[s,t] \oplus (p_2r_3 - p_3r_2)\mathbb{C}[s,t] \oplus pw(p_3q_2-p_2q_1)\mathbb{C}[s,t] \\
        &v_3 \in -p(q_2r_1+q_1r_2)\mathbb{C}[s,t] \oplus (p_2r_1 - p_1r_2)\mathbb{C}[s,t] \oplus pw(p_2q_1+p_1q_2)\mathbb{C}[s,t].
    \end{align*}
    where
    \begin{equation*}
        \begin{bmatrix}
            p_1 \\
            p_2 \\
            p_3
        \end{bmatrix}
        =
        \begin{bmatrix}
            2xy \\
            y^2-x^2 \\
            y^2+x^2
        \end{bmatrix}, \:
        \begin{bmatrix}
            q_1 \\
            q_2 \\
            q_3
        \end{bmatrix}
        =
        \begin{bmatrix}
            p_2p_3 \\
            -p_1p_3 \\
            p_1p_2
        \end{bmatrix}, \:
        \begin{bmatrix}
            r_1 \\
            r_2 \\
            r_3
        \end{bmatrix}
        =
        \begin{bmatrix}
            p_1^3 \\
            p_2^3 \\
            -p_3^3
        \end{bmatrix},
    \end{equation*}
    and \(p = p_1p_2p_3 = 2xy(y^4-x^4)\) and \(w = x^2y^2(x^4-y^4)^2\) are as in Lemma \ref{semi-invGI}. 
\end{theorem}

\begin{proof}
    We saw in Lemma \ref{W_L:n=0(8)} that \(v_0\) is an invariant and hence lies in \(\mathbb{C}[s,t]\), implying that \(n = \deg v_0 \equiv 0 \ (\textrm{mod}\ 8)\). Assume first that all odd weights are congruent to \(1\) modulo \(4\). We saw in the proof of Lemma \ref{semi-invGI} that 
    \begin{align*}
        &A\cdot p_1 = -p_2,\, A\cdot p_2 = p_1,\, A\cdot p_3 = p_3 \\
        &B\cdot p_1 = -ip_1,\, B\cdot p_2 = p_3,\, B\cdot p_3 = p_2.
    \end{align*}
    It follows that \(p_2v_1 - p_1v_2 + p_3v_3\) is a semi-invariant of weight \(\chi_i\) (change \(i\) accordingly as in table) and, hence, belongs to \(pu\mathbb{C}[s,t]\). It also follows that \(G\) acts on each \(q_i\) and each \(r_i\) in the following way:
    \begin{align*}
        &A\cdot q_1 = -q_2,\, A\cdot q_2 = q_1,\, A\cdot q_3 = -q_3 \\
        &B\cdot q_1 = q_1,\, B\cdot q_2 = iq_3,\, B\cdot q_3 = iq_2 
    \end{align*}
    and
    \begin{align*}
        &A\cdot r_1 = -r_2,\, A\cdot r_2 = r_1,\, A\cdot r_3 = r_3 \\
        &B\cdot r_1 = ir_1,\, B\cdot r_2 = -r_3,\, B\cdot r_3 = -r_2. 
    \end{align*}
    This, along with Lemmas \ref{repG_n=0(4)} amd \ref{semi-invGI}, implies that there are invariants \(f, g, h \in \mathbb{C}[s,t]\) such that 
    \begin{equation*}
        \begin{bmatrix}
            p_2 & -p_1 & p_3 \\
            q_2 & q_1 & q_3 \\
            r_2 & -r_1 & r_3
        \end{bmatrix}
        \begin{bmatrix}
            v_1 \\
            v_2 \\
            v_3
        \end{bmatrix}
        =
        \begin{bmatrix}
            (pu)f \\
            ug \\
            (puw)h
        \end{bmatrix}
    \end{equation*}
    where \(u\) is as in Lemma \ref{semi-invGI}. Letting \(N\) be the coefficient matrix, it can be seen that \(\det N = 4u\), giving
    \begin{equation*}
        \begin{bmatrix}
            v_1 \\
            v_2 \\
            v_3
        \end{bmatrix}
        = \frac{1}{4}M
        \begin{bmatrix}
            pf \\
            g \\
            pwh
        \end{bmatrix} =
        \frac{1}{4} \tilde{M}
        \begin{bmatrix}
            f \\
            g \\ 
            h
        \end{bmatrix}
    \end{equation*}
    where
    \begin{equation*}
        M = 
        \begin{bmatrix}
            q_1r_3+q_3r_1 & p_1r_3 - p_3r_1 & -(p_1q_3+p_3q_1) \\
            q_3r_2-q_2r_3 & p_2r_3 - p_3r_2 & p_3q_2-p_2q_1 \\
            -(q_2r_1+q_1r_2) & p_2r_1 - p_1r_2 & p_2q_1+p_1q_2
        \end{bmatrix}
    \end{equation*}
    is the adjugate of \(N\) and
    \begin{equation*}
        \tilde{M} = M
        \begin{bmatrix}
            p & 0 & 0 \\
            0 & 1 & 0 \\
            0 & 0 & pw
        \end{bmatrix},
    \end{equation*}
    so the \(i\)th row of \(\tilde{M}\) consists of the generators of the module that \(v_i\) is claimed to belong to. It remains to check that these generators are free.
    
    Since \(A \cdot \tilde{M}_{1j} = -\tilde{M}_{2j}\) and \(B \cdot \tilde{M}_{1j} = -i\tilde{M}_{3j}\) for \(j = 1,2,3\), it suffices to show that \(\tilde{M}_{11},\, \tilde{M}_{12}\), and \(\tilde{M}_{13}\) are linearly independent over \(\mathbb{C}[s,t]\). If \(f_1, f_2, f_3 \in \mathbb{C}[s,t]\) are such that \(\tilde{M}_{11}f_1 + \tilde{M}_{12}f_2 + \tilde{M}_{13}f_3 = 0 \), then the action of \(A\) and \(B\) on the first row of \(\tilde{M}\) implies that 
    \begin{equation*}
        \tilde{M}
        \begin{bmatrix}
            f_1 \\ 
            f_2 \\ 
            f_3
        \end{bmatrix}
        = 0.
    \end{equation*}
    Since \(\det \tilde{M} = p^2w \det M = (2w)^2 (\det N)^2 = (8wu)^2 \neq 0\), we conclude that \(f_i = 0\).
    
    For the case when all odd weights are congruent to \(3\) modulo \(4\), simply replace \(p_3\) by \(-p_3\) and apply the same argument.
\end{proof}

\end{document}
